% \iffalse meta-comment
%
% hit-math.dtx 是各个类共用的数学层
%
% \fi
%
% \section{共享数学层}
%
% \changes{v3.2a}{2026/08/07}{数学部分从各个类的 \texttt{*-deps-a} 模块抽出合并}
% \pkg{amsmath} 一路到数学字体选择这一段，原来每个类一字不差地各写了一遍。
% 这里合成一份。
%
% 位置与 \file{hit-fonts.dtx} 一样，排在类文件头和 \file{hit-fonts.dtx} 之后、类
% 自己的 dtx 之前。这个位置在 \cs{LoadClass} 之前，宏包还不能加载，所以只留声明，
% 其余放进 \cs{hit@setup@math}，由各个类在原来的位置调用。
%
%    \begin{macrocode}
%<*shared-math>
\prop_new:N \g__hit_theorem_option_prop
\seq_new:N \g__hit_theorem_order_seq
\tl_new:N \g__hit_theorem_within_tl
\tl_new:N \g__hit_theorem_option_list_tl
\tl_new:N \l__hit_theorem_temporary_tl
\int_new:N \g__hit_theorem_style_int
\tl_new:N \g__hit_theorem_style_tl
\str_new:N \l__hit_math_font_str
\str_new:N \l__hit_math_style_str
\clist_new:N \l__hit_math_options_clist
\bool_new:N \g__hit_unicodemath_bool
\cs_new_protected:Npn \hit@setup@math
  {
%    \end{macrocode}
%    \begin{macrocode}
\RequirePackage{amsmath}
%    \end{macrocode}
%
% 定理类环境宏包，其中 \pkg{amsmath} 选项用来兼容 \AmSTeX\ 的宏包
%    \begin{macrocode}
\RequirePackage{amsthm}
\RequirePackage{thmtools}
%    \end{macrocode}
%
% \changes{v3.2a}{2026/08/06}{补回 \pkg{aliasctr} 覆盖 \cs{@addtoreset} 时丢掉的
%   \cs{theH} 定义（\issue[183]）}
% \pkg{thmtools} 依赖的 \pkg{aliasctr} 用 \cs{renewcommand*} 平覆盖了内核的
% \cs{@addtoreset}，只保留计数器清零链，丢掉了内核同时生成的 \cs{theH} 定义。
% LaTeX 2024-11-01 起 \pkg{hyperref} 不再自己打这个补丁（\cs{hyper@nopatch@counter}
% 恒定义），于是 \pkg{subfigure}、\pkg{lstlisting} 和定理环境的 PDF 锚点都丢了章号
% 前缀。这里在保留 \pkg{aliasctr} 别名跟随行为的前提下把 \cs{theH} 补回去。
%    \begin{macrocode}
\cs_set_eq:NN \hit@original@add@to@reset \@addtoreset
\cs_set:Npn \@addtoreset ##1##2
  {
    \hit@original@add@to@reset {##1} {##2}
    \cs_gset:cpx { theH ##1 }
      { \exp_not:c { theH ##2 } . \exp_not:N \the \exp_not:N \value {##1} }
  }
%    \end{macrocode}
%
% \changes{v3.2a}{2026/08/07}{定理环境的样式与声明收进 \cs{hit@theorem@style} 与
%   \cs{hit@theorem@list}，为 \issue[183] 换 \pkg{thmtools} 做准备}
% 定理环境的内部接口。样式与声明分散在三处（全局样式、\file{cfg} 里的
% \env{proof}、十几个编号环境），换实现时要改的地方太多，先收到这里。
% \cs{hit@theorem@style} 的键名按 \pkg{thmtools} 的叫法取，将来换实现只改这一处的
% 映射；\cs{hit@theorem@list} 接一张「环境名 = 标题」的表，顺序即声明顺序。
% \cs{hit@theorem@declare@star@list} 是不编号那批的同款入口。
%
% 三个都从 \cs{hitsetup} 的 \option{thm} 族有键可达（样式键、\option{list}、
% \option{star}），配置文件因此不必直接调命令。
%    \begin{macrocode}
\cs_new_protected:Npn \__hit_theorem_option:nn ##1##2
  {
    \prop_if_in:NnF \g__hit_theorem_option_prop {##1}
      { \seq_gput_right:Nn \g__hit_theorem_order_seq {##1} }
    \prop_gput:Nnn \g__hit_theorem_option_prop {##1} {##2}
  }
\cs_new_protected:Npn \__hit_theorem_build:
  {
    \tl_gclear:N \g__hit_theorem_option_list_tl
    \seq_map_inline:Nn \g__hit_theorem_order_seq
      {
        \prop_get:NnN \g__hit_theorem_option_prop {####1} \l__hit_theorem_temporary_tl
        \tl_gput_right:Nx \g__hit_theorem_option_list_tl
          { ####1 = { \exp_not:V \l__hit_theorem_temporary_tl } , }
      }
  }
%    \end{macrocode}
%
% \pkg{thmtools} 的样式键只认 \cs{declaretheoremstyle}，\cs{declaretheorem} 不吃。
% 所以每次改样式就生成一个新样式名，声明环境时引用它。名字带序号是因为
% \pkg{amsthm} 的 \cs{newtheoremstyle} 重名会报错。
%    \begin{macrocode}
\cs_new_protected:Npn \__hit_theorem_new_style:
  {
    \__hit_theorem_build:
    \int_gincr:N \g__hit_theorem_style_int
    \tl_gset:Nx \g__hit_theorem_style_tl { hitthm \int_use:N \g__hit_theorem_style_int }
    \use:x
      {
        \exp_not:N \declaretheoremstyle
          [ \exp_not:V \g__hit_theorem_option_list_tl ] { \g__hit_theorem_style_tl }
      }
    \use:x { \exp_not:N \theoremstyle { \g__hit_theorem_style_tl } }
  }
\cs_new_protected:Npn \hit@theorem@set@auto@ref@name ##1##2
  { \cs_gset:cpn { ##1 autorefname } {##2} }
\cs_new_protected:Npn \hit@theorem@declare ##1##2
  {
    % \pkg{amsthm} 自带 proof，重名的先撤掉再声明
    \cs_undefine:c {##1} \cs_undefine:c { end ##1 }
    \use:x
      {
        \exp_not:N \declaretheorem
          [
            style = \g__hit_theorem_style_tl ,
            name = { \exp_not:n {##2} } ,
            \tl_if_empty:NF \g__hit_theorem_within_tl
              { numberwithin = \exp_not:V \g__hit_theorem_within_tl , }
          ]
          {##1}
      }
    \hit@theorem@set@auto@ref@name {##1} {##2}
  }
\cs_new_protected:Npn \hit@theorem@declare@star ##1##2
  {
    \cs_undefine:c {##1} \cs_undefine:c { end ##1 }
    \use:x
      {
        \exp_not:N \declaretheorem
          [ style = \g__hit_theorem_style_tl , name = { \exp_not:n {##2} } , numbered = no ]
          {##1}
      }
    \hit@theorem@set@auto@ref@name {##1} {##2}
  }
\keys_define:nn { hit / thm }
  {
    body-font   .code:n = { \__hit_theorem_option:nn { bodyfont }   {##1} } ,
    head-font   .code:n = { \__hit_theorem_option:nn { headfont }   {##1} } ,
    qed         .code:n = { \__hit_theorem_option:nn { qed }        {##1} } ,
    head-punct  .code:n = { \__hit_theorem_option:nn { headpunct }  {##1} } ,
    note-braces .code:n = { \__hit_theorem_option:nn { notebraces } {##1} } ,
    note-font   .code:n = { \__hit_theorem_option:nn { notefont }   {##1} } ,
    space-above .code:n = { \__hit_theorem_option:nn { spaceabove } {##1} } ,
    space-below .code:n = { \__hit_theorem_option:nn { spacebelow } {##1} } ,
    head-space  .code:n = { \__hit_theorem_option:nn { postheadspace } {##1} } ,
    head-indent .code:n = { \__hit_theorem_option:nn { headindent } {##1} } ,
    within .code:n = { \tl_gset:Nn \g__hit_theorem_within_tl {##1} } ,
    list .code:n = { \hit@theorem@list {##1} } ,
    star .code:n = { \hit@theorem@declare@star@list {##1} } ,
  }
\cs_new_protected:Npn \hit@theorem@style ##1
  { \hit@keys@set:nn { hit / thm } {##1} \__hit_theorem_new_style: }
\cs_new_protected:Npn \hit@theorem@list ##1
  { \keyval_parse:nnn { \use_none:n } { \hit@theorem@declare } {##1} }
\cs_new_protected:Npn \hit@theorem@declare@star@list ##1
  { \keyval_parse:nnn { \use_none:n } { \hit@theorem@declare@star } {##1} }
%    \end{macrocode}
%
% \changes{v3.2a}{2026/08/06}{为 \pkg{ntheorem} 专有命令加兼容层（\issue[183]）}
% \pkg{ntheorem} 专有命令的兼容层。前四条能一一对应到 \cs{hitsetup} 的
% \option{theorem} 键，转发过去；后面几条 \pkg{thmtools} 没有对应写法，报个警告
% 后忽略。\cs{theoremstyle} 不动，那是 \pkg{amsthm} 自己的命令。
%    \begin{macrocode}
\cs_new_protected:Npn \__hit_theorem_legacy:nnn ##1##2##3
  {
    \__hit_warn_once:nnn {##1} { deprecated-ntheorem } {##2}
    \hit@theorem@style { ##2 = {##3} }
  }
\cs_new_protected:Npn \__hit_theorem_dropped:nn ##1##2
  { \__hit_warn_once:nnn {##1} { removed-ntheorem } {##2} }
\cs_new_protected:Npn \theorembodyfont ##1
  { \__hit_theorem_legacy:nnn { theorembodyfont } { body-font } {##1} }
\cs_new_protected:Npn \theoremheaderfont ##1
  { \__hit_theorem_legacy:nnn { theoremheaderfont } { head-font } {##1} }
\cs_new_protected:Npn \theoremsymbol ##1
  { \__hit_theorem_legacy:nnn { theoremsymbol } { qed } {##1} }
\cs_new_protected:Npn \theoremseparator ##1
  { \__hit_theorem_legacy:nnn { theoremseparator } { head-punct } {##1} }
\cs_new_protected:Npn \theoremnumbering ##1
  { \__hit_warn_once:nnn { theoremnumbering } { removed-ntheorem-numbering } { } }
\cs_new_protected:Npn \theoremprework ##1
  { \__hit_theorem_dropped:nn { theoremprework } { declaretheorem } }
\cs_new_protected:Npn \theorempostwork ##1
  { \__hit_theorem_dropped:nn { theorempostwork } { declaretheorem } }
\cs_new_protected:Npn \theoremclass ##1
  { \__hit_theorem_dropped:nn { theoremclass } { declaretheorem } }
\cs_new_protected:Npn \theoremlisttype ##1
  { \__hit_theorem_dropped:nn { theoremlisttype } { listoftheorems } }
\cs_new_protected:Npn \listtheorems ##1
  {
    \__hit_theorem_dropped:nn { listtheorems } { listoftheorems }
    \listoftheorems [ ignoreall , show = {##1} ]
  }
\cs_new_protected:Npn \thref ##1
  {
    \__hit_theorem_dropped:nn { thref } { autoref }
    \autoref {##1}
  }
%    \end{macrocode}
%
% 删除 \pkg{newtxtext}，由于最新一次更新 (\texttt{revision 62369}) 中与 \pkg{ctex} 宏包冲突，且无法解决，将正文字体替换为 TG TermesX 与 TG Heros.
% \changes{v3.0.18}{2022/05/02}{删除 \pkg{newtxtext}，由于最新一次更新 (\texttt{revision 62369}) 中与 \pkg{ctex} 宏包冲突，且无法解决，将正文字体替换为 TG Termes 与 TG Heros}
% \changes{v3.0.19}{2022/05/05}{将 TG Termes 字体替换为 TG TermesX 字体，与 \pkg{newtxmath} 宏包更适配，删除 \option{Scale} 选项}
% \changes{v3.2a}{2026/08/06}{英文字体三件套设置从 bookcls 拆出至 book-fonts 模块（与 art-fonts / artplus-fonts 保持一致）}
%
% 单独使用 \pkg{newtxmath} 宏包改变数学字体时，会使 \texttt{operators}，\cs{mathrm}，\cs{mathit}，\cs{mathbf} 命令使用 CM 字体，根据 \href{https://github.com/sjtug/SJTUThesis/blob/b164d0f5b3087ed86c1038b5a85014b47a92c6f2/texmf/tex/latex/sjtuthesis/fd/sjtu-math-font-termes.def#L52-L65}{SJTUThesis 的字体设置} 进行调整
% \changes{v3.0.19}{2022/05/05}{对 \pkg{newtxmath} 宏包打补丁，修复 \texttt{operators}，\cs{mathrm}，\cs{mathit}，\cs{mathbf} 命令的字体问题}
%    \begin{macrocode}
% \changes{v3.2a}{2026/08/15}{\option{fontset-math} 的 \texttt{auto} 默认从
%   \texttt{xits} 改 \texttt{termes}}
% \texttt{auto} 走 \texttt{termes}（TeX Gyre Termes Math）而不是 \texttt{xits}：
% 学校要的正文西文是 Times，Termes 是 Times 的克隆，数学与正文同一血统；
% XITS 源自 STIX，空心体那批字形是几何化的圆头，跟 Times 摆在一起不像一家。
% 代价是覆盖面窄一些——Termes 1112 个码位、XITS 2423 个，缺口在冷门运算符、
% 箭头与杂项技术符号（2300 段 51 对 90）。常用符号两边都全。
% 要覆盖面就写 \option{fontset-math}\texttt{=stix}，STIX~Two 有 3607 个码位，
% 比 XITS 还多。
%
% \changes{v3.2a}{2026/08/15}{\option{math-style} 默认改 \texttt{auto}，
%   显式写了又走 \texttt{newtx} 时提示}
% \option{math-style} 默认 \texttt{auto}，行为与从前的 \texttt{GB} 一致。
% 改成 \texttt{auto} 是为了能分辨「用户显式写过」——那组选项只对
% \pkg{unicode-math} 那几档有效，\texttt{newtx} 走 8~位路径，希腊字母正斜体由
% \pkg{newtxmath} 自己决定。两者撞上时发一条提示，默认值撞上则不打扰。
\str_set:Nx \l__hit_math_font_str { \hit@mathfont }
\str_if_eq:VnT \l__hit_math_font_str { } { \str_set:Nn \l__hit_math_font_str { auto } }
\str_if_eq:VnT \l__hit_math_font_str { auto }
  {
    \bool_if:NTF \g__hit_newtxmath_bool
      { \str_set:Nn \l__hit_math_font_str { newtx } }
      { \str_set:Nn \l__hit_math_font_str { termes } }
  }
%    \end{macrocode}
%
% 数学排版风格。\texttt{TeX} 是 \hologo{TeX} 传统写法（大写希腊字母直立、小写倾斜），
% \texttt{ISO} 按 ISO 80000-2（希腊字母一律倾斜、粗体同理），\texttt{GB} 在 ISO
% 的基础上把偏导符号 \cs{partial} 改直立，这是国标的要求。
%
% 五个 \texttt{math*=sym} 是为了照顾旧文档。\pkg{unicode-math} 默认让
% \cs{mathrm}、\cs{mathbf} 这几个走正文字体，正文字体里没有数学希腊字母，
% \cs{mathbf} 套希腊字母这种写法会静悄悄掉字；改成 \texttt{sym} 就跟原先的
% 8~位数学字体一样从数学字体取字。
%    \begin{macrocode}
\str_set:Nx \l__hit_math_style_str { \hit@mathstyle }
\bool_new:N \l__hit_math_style_explicit_bool
\str_if_eq:VnTF \l__hit_math_style_str { }
  { \str_set:Nn \l__hit_math_style_str { auto } }
  { \bool_set_true:N \l__hit_math_style_explicit_bool }
\str_if_eq:VnT \l__hit_math_style_str { auto }
  {
    \bool_set_false:N \l__hit_math_style_explicit_bool
    \str_set:Nn \l__hit_math_style_str { GB }
  }
\bool_lazy_and:nnT
  { \bool_if_p:N \l__hit_math_style_explicit_bool }
  { \str_if_eq_p:Vn \l__hit_math_font_str { newtx } }
  { \msg_warning:nn { hithesis } { math-style-with-newtx } }
\str_case:VnF \l__hit_math_style_str
  {
    { TeX } { \clist_set:Nn \l__hit_math_options_clist { math-style=TeX , bold-style=TeX , mathrm=sym , mathit=sym , mathbf=sym , mathsf=sym , mathtt=sym } }
    { ISO } { \clist_set:Nn \l__hit_math_options_clist { math-style=ISO , bold-style=ISO , mathrm=sym , mathit=sym , mathbf=sym , mathsf=sym , mathtt=sym } }
    { GB }
      { \clist_set:Nn \l__hit_math_options_clist
          { math-style=ISO , bold-style=ISO , partial=upright , mathrm=sym , mathit=sym , mathbf=sym , mathsf=sym , mathtt=sym } }
  }
  { \msg_error:nnx { hithesis } { math-style-unknown } { \l__hit_math_style_str } }
%    \end{macrocode}
%
% \pkg{unicode-math} 各档对应的字体文件。第二个参数是粗体数学版本，
% \pkg{amsmath} 的 \cs{boldsymbol} 要用它；没有独立粗体的字体就填回常规体，
% 不填会报 \texttt{Font shape .../b/n undefined} 并掉字。\texttt{newtx} 与
% \texttt{none} 不走 \pkg{unicode-math}，单列在后面。
%    \begin{macrocode}
\cs_new_protected:Npn \__hit_set_math_font:nn ##1##2
  {
    \bool_gset_true:N \g__hit_unicodemath_bool
    \exp_args:NV \PassOptionsToPackage \l__hit_math_options_clist { unicode-math }
    \RequirePackage { unicode-math }
    \setmathfont {##1}
    \setmathfont {##2} [ version = bold ]
%    \end{macrocode}
%
% \pkg{unicode-math} 的 \cs{symbf} 跟着 \texttt{bold-style} 走，ISO 下是粗斜体，
% 而 \cs{mathbf} 在传统 \hologo{TeX} 里一直是粗正体。把 \cs{mathbf} 单独接到
% \cs{symbfup} 上，旧文档里的 \cs{mathbf} 才不会变成斜体。\pkg{unicode-math}
% 在 \cs{begin\{document\}} 时会重新安排一遍数学字母表，所以要挂在钩子里。
%    \begin{macrocode}
    \AddToHook { begindocument/end } { \cs_set_eq:NN \mathbf \symbfup }
%    \end{macrocode}
%
% \changes{v3.2a}{2026/08/10}{\pkg{unicode-math} 路径上加一道 \pkg{bm} 的补救钩子。
%   v3.1e 的 \file{hithesis.sty} 无条件 \cs{RequirePackage}\marg{bm}，用户自己改过的
%   那份升级后照旧优先加载，完整示例实测 70 个 Extended mathchar used as mathchar}
% \pkg{bm} 跟 \pkg{unicode-math} 不兼容：它把 \cs{boldsymbol} 换成只认 8~位
% mathchar 的版本，公式里一出现希腊字母就报 \texttt{Extended mathchar used as
% mathchar}。模板自己不装 \pkg{bm}，但用户的导言区或者自己那份 \file{hithesis.sty}
% 会装，所以挂个钩子：它真被装上时，把 \cs{bm} 与 \cs{boldsymbol} 接回
% \pkg{unicode-math} 自己的 \cs{symbf}，并说明一句。\pkg{bm} 没被装就什么也不做。
%    \begin{macrocode}
    \AddToHook { package/bm/after }
      {
        \cs_set_eq:NN \bm \symbf
        \cs_set_eq:NN \boldsymbol \symbf
        \msg_warning:nn { hithesis } { bm-with-unicode-math }
      }
  }
\cs_new_protected:Npn \__hit_set_math_xits:
  { \__hit_set_math_font:nn { XITSMath-Regular.otf } { XITSMath-Bold.otf } }
\cs_new_protected:Npn \__hit_set_math_stix:
  { \__hit_set_math_font:nn { STIXTwoMath-Regular.otf } { STIXTwoMath-Regular.otf } }
\cs_new_protected:Npn \__hit_set_math_termes:
  { \__hit_set_math_font:nn { texgyretermes-math.otf } { texgyretermes-math.otf } }
\cs_new_protected:Npn \__hit_set_math_libertinus:
  { \__hit_set_math_font:nn { libertinusmath-regular.otf } { libertinusmath-bold.otf } }
\cs_new_protected:Npn \__hit_set_math_newcm:
  { \__hit_set_math_font:nn { NewCMMath-Regular.otf } { NewCMMath-Bold.otf } }
\cs_new_protected:Npn \__hit_set_math_latinmodern:
  { \__hit_set_math_font:nn { latinmodern-math.otf } { latinmodern-math.otf } }
%    \end{macrocode}
%
% \changes{v3.1b}{2022/06/06}{直接修改编码，而不引入新的宏包}
% \changes{v3.2a}{2026/08/08}{newtxmath 的加载分支迁到 expl3}
% \pkg{newtxmath} 是 8~位字体路径，加载时得临时把三件套换成 \pkg{newtx} 自己的
% 编码与字族，否则 \texttt{operators}、\cs{mathrm}、\cs{mathit}、\cs{mathbf}
% 会退回 CM。做法取自
% \href{https://github.com/sjtug/SJTUThesis/blob/b164d0f5b3087ed86c1038b5a85014b47a92c6f2/texmf/tex/latex/sjtuthesis/fd/sjtu-math-font-termes.def#L52-L65}{SJTUThesis 的字体设置}。
%
% \changes{v3.2a}{2026/08/15}{\texttt{newtx} 那档补上 \cs{symbf} 一组的等价写法}
% \pkg{unicode-math} 的 \cs{symbf}、\cs{symup}、\cs{symit} 在 8~位路径上不存在，
% 同一份文档从别的档切到 \texttt{newtx}，正文里每一处 \cs{symbf} 都会报未定义。
% 这里接出三个等价写法，让文档在两条路径上都编得过：
% \begin{center}
% \begin{tabular}{ll}
% \toprule
% \pkg{unicode-math} & \texttt{newtx} 这边接到 \\\midrule
% \cs{symbf} & \cs{bm} \\
% \cs{symup} & \cs{mathrm} \\
% \cs{symit} & \cs{mathit} \\
% \bottomrule
% \end{tabular}
% \end{center}
%
% 不是逐字等价：\cs{symbf} 在 \pkg{unicode-math} 那边跟着 \option{bold-style}
% 走（ISO 下是粗斜体），\cs{bm} 不跟。要的是「切过去能编、观感接近」，
% 不是两条路径排出同一个 PDF——那本来也做不到，两套字体的字形不一样。
%
% 接在 \texttt{begindocument/end} 上：\cs{bm} 由用户的 \cs{usepackage}\marg{bm}
% 提供，类加载时还不存在，那时候接等于把 \cs{symbf} 定义成一个未定义的东西。
% 用户没装 \pkg{bm} 时 \cs{symbf} 退到 \cs{mathbf}——粗体是有的，只是不跟着
% \option{bold-style} 走。已经有定义就不动，用户自己定义过的优先。
%    \begin{macrocode}
\cs_new_protected:Npn \__hit_math_sym_shim:
  {
    \AddToHook { begindocument/end }
      {
        \cs_if_exist:NF \symbf
          {
            \cs_if_exist:NTF \bm
              { \cs_new_eq:NN \symbf \bm }
              { \cs_new_eq:NN \symbf \mathbf }
          }
        \cs_if_exist:NF \symup { \cs_new_eq:NN \symup \mathrm }
        \cs_if_exist:NF \symit { \cs_new_eq:NN \symit \mathit }
      }
  }
%    \end{macrocode}
%
%    \begin{macrocode}
\cs_new_protected:Npn \__hit_set_math_newtx:
  {
    \cs_set_eq:NN \oldencodingdefault \encodingdefault
    \cs_set_eq:NN \oldrmdefault \rmdefault
    \cs_set_eq:NN \oldsfdefault \sfdefault
    \cs_set_eq:NN \oldttdefault \ttdefault
    \tl_set:Nn \encodingdefault { T1 }
    \tl_set:Nn \rmdefault { ntxtlf }
    \tl_set:Nn \sfdefault { qhv }
    \tl_set:Nn \ttdefault { ntxtt }
    \RequirePackage { newtxmath }
    \cs_set_eq:NN \encodingdefault \oldencodingdefault
    \cs_set_eq:NN \rmdefault \oldrmdefault
    \cs_set_eq:NN \sfdefault \oldsfdefault
    \cs_set_eq:NN \ttdefault \oldttdefault
    \__hit_math_sym_shim:
  }
\cs_new_protected:Npn \__hit_set_math_none: { }
%    \end{macrocode}
%
% \pkg{amssymb} 与 \pkg{unicode-math} 冲突，只在 8~位路径上加载；走
% \pkg{unicode-math} 时它提供的符号由数学字体自带。它要排在 \pkg{newtxmath}
% 前面，两者都定义 \cs{Bbbk}，反过来会报重复定义。
%    \begin{macrocode}
\cs_if_exist:cTF { __hit_set_math_ \l__hit_math_font_str : }
  {
    \bool_lazy_or:nnT
      { \str_if_eq_p:Vn \l__hit_math_font_str { newtx } }
      { \str_if_eq_p:Vn \l__hit_math_font_str { none } }
      { \RequirePackage { amssymb } }
    \use:c { __hit_set_math_ \l__hit_math_font_str : }
  }
  { \msg_error:nnx { hithesis } { math-font-unknown } { \l__hit_math_font_str } }
  }
%    \end{macrocode}
%
% \changes{v3.2a}{2026/08/06}{新增 \env{eqdenote} 环境，公式的物理量注释}
% \env{eqdenote} 排公式下面那张「式中……」注释表。规范要求注释的转行与破折号后
% 第一个字对齐，所以底下是一个 \pkg{tabularx}，说明列取 \texttt{X}。
%
% 排出来是三列：引导词、符号、说明。用户只写后两列，第二行起那个空的引导词格子
% 由环境自己补。可选参数换引导词，传空则留出等宽的空位，拆成两段接排时左缘对得上。
%    \begin{macrocode}
\dim_new:N \l__hit_eqdenote_depth_dim
\box_new:N \l__hit_eqdenote_box
\seq_new:N \l__hit_eqdenote_seq
\tl_new:N \l__hit_eqdenote_tl
\cs_new:Npn \hit@eqdenote@label
  {
    \bool_if_exist:NTF \g__hit_en_bool
      {
        \bool_if:NTF \g__hit_en_bool
          { \hit@eqdenote@label@en } { \hit@eqdenote@label@zh }
      }
      { \hit@eqdenote@label@zh }
  }
%    \end{macrocode}
%
% 表格要用 \cs{begin}\marg{tabularx} 与 \cs{end}\marg{tabularx} 成对写在一条命令
% 里。\pkg{tabularx} 会把正文反复排好几遍来定 \texttt{X} 列宽度，直接在环境的开始
% 代码里写 \cs{tabularx}，收尾那几句也会被卷进去重排，\cs{prevdepth} 在受限水平
% 模式下会报错。
%    \begin{macrocode}
\tl_new:N \l__hit_eqdenote_font_tl
%    \end{macrocode}
%
% \pkg{longtable} 被模板改成了五号（\cs{LT@array} 里加了 \cs{wuhao}），那是给插表
% 用的。式中注释是正文的一部分，字号要跟正文一样，所以进表格前用 \cs{ltfontsize}
% 把当前字号还回去。字号对了，\cs{ccwd}、\cs{strutbox} 与行距也就都跟正文一致。
%    \begin{macrocode}
\cs_new_protected:Npn \__hit_eqdenote_keep_font:
  {
    \tl_set:Nx \l__hit_eqdenote_font_tl
      {
        \exp_not:N \fontsize { \f@size } { \f@baselineskip }
        \exp_not:N \selectfont
      }
    \exp_args:NV \ltfontsize \l__hit_eqdenote_font_tl
  }
%    \end{macrocode}
%
% \changes{v3.2a}{2026/08/06}{\env{eqdenote} 底下从 \pkg{tabularx} 换成
%   \pkg{xltabular}，注释表放不下时自动转页接排，不必手动拆成两个环境}
% \pkg{xltabular} 是 \pkg{longtable} 加 \texttt{X} 列，表格能自动跨页接排。
% \cs{LTleft} 与 \cs{LTright} 让表格靠左，\cs{LTpost} 去掉 \pkg{longtable} 默认在
% 表格后面加的 \cs{bigskipamount}。
%
% \changes{v3.2a}{2026/08/07}{\env{eqdenote} 的段前距离补上 \pkg{longtable} 漏掉的
%   行距胶，之前接正文比正常行距紧 4.4pt}
% \cs{LTpre} 不能简单置零。\pkg{longtable} 把首块排在一个新 vbox 里
%（\cs{LT@bchunk} 的 \cs{setbox}\cs{z@}\cs{vbox}），vbox 开头 \cs{prevdepth} 是
% \texttt{-1000pt}，TeX 不给首行插行距胶；\cs{unvbox} 出来时也不补。于是表首行的基线
% 与上一行正文之间只隔「上一行的深度加首行的高度」，比正常行距紧。量出来是 4.4278pt，
% 而 \cs{baselineskip} 减去这两项正好是这个数，可见缺的就是那份胶。
%
% \cs{\_\_hit\_eqdenote\_pre\_skip:} 把它算出来放进 \cs{LTpre}。先 \cs{par}，
% 否则 \cs{prevdepth} 还是上一段的。算出来比 \cs{lineskiplimit} 小时改用
% \cs{lineskip}，与 TeX 自己的规则一致。\cs{prevdepth} 是 \texttt{-1000pt}（页首、
% 紧跟标题）时本来就不该有胶，置零。
%    \begin{macrocode}
%    \end{macrocode}
%
% \changes{v3.2a}{2026/08/18}{\env{eqdenote} 关段时停掉段末撑竿钩子：公式后
%   的空续段被撑竿撑成一行幽灵行，段前多一个行距，断页时还抢走页首基线}
% \emph{关段要避开撑竿钩子。} \cs{end}\marg{equation} 之后段落还开着（行间
% 公式属于所在段落），紧跟 \env{eqdenote} 时这个续段是空的。\hologo{TeX} 对
% 空续段的 \cs{par} 本不产生任何行，但 \texttt{para/end} 钩子里的段落撑竿把
% 它撑成一行只有零宽 rule 的幽灵行：行中形态下公式到「式中」多出一个行距
% （实测 48 对范例 22），表被推到新页时幽灵行抢走 \cs{topskip} 基线，
% 「式中」首行比正文页首低一行（实测 140.5 对 119.9）。关段那一下把
% \cs{hit@format@strut:} 临时置空，空续段不再成行；前一行真正文的撑竿在
% 各自成行时已经排进去了，不受影响。
%    \begin{macrocode}
\cs_new_protected:Npn \__hit_eqdenote_pre_skip:
  {
    \group_begin:
    \cs_if_exist:NT \hit@format@strut:
      { \cs_set_eq:NN \hit@format@strut: \prg_do_nothing: }
    \par
    \group_end:
    \dim_compare:nNnTF { \tex_prevdepth:D } < { -999pt }
      { \skip_zero:N \LTpre }
      {
        \skip_set:Nn \LTpre
          { \baselineskip - \tex_prevdepth:D - \arraystretch \ht \strutbox }
        \dim_compare:nNnT { \LTpre } < { \lineskiplimit }
          { \skip_set_eq:NN \LTpre \lineskip }
      }
  }
\cs_new_protected:Npn \__hit_eqdenote_table:nn #1#2
  {
    \__hit_eqdenote_keep_font:
%    \end{macrocode}
%
% \changes{v3.2a}{2026/08/18}{\env{eqdenote} 的行撑竿缩到 0.9 倍：撑竿恰好等于
%   行距时 \hologo{TeX} 退到 lineskip 堆叠，行距多出约 1\,pt，换页后首行基线也
%   不落在 \cs{topskip} 上}
% 行撑竿要\emph{小于}行距。\cs{strutbox} 的高深加起来正好一个 \cs{baselineskip}，
% \cs{arraystretch} 取 1 时表格行高不小于行距，基线间距放不下，\hologo{TeX} 退到
% 按 \cs{lineskip} 堆叠，行距变成「行高加 1\,pt」（实测 20.9 对正文 19.45）；
% 换页后首行的基线取行高而不是 \cs{topskip}，比正文页首低约 22\,bp——1--7 页版
% 实测正文页首基线 119.9，注释表首行 141.8。缩到 0.89 后撑竿高 12.16 低于
% 正文撑竿 12.22，行盒尺寸由格内段落撑竿接管，恰为一个 \cs{baselineskip}，
% 行距与页首基线两处一起归位（0.9 时撑竿高 12.30 仍比上限多 0.08\,pt，
% 还是掉进 lineskip 堆叠）。
%    \begin{macrocode}
    \cs_set:Npn \arraystretch { 0.89 }
    \skip_zero:N \LTleft
    \skip_set:Nn \LTright { 0pt~plus~1fill }
    \__hit_eqdenote_pre_skip:
    \skip_zero:N \LTpost
%    \end{macrocode}
%
% \changes{v3.2a}{2026/08/18}{清空 \pkg{longtable} 的表头盒。零高的空表头盒
%   恰好落在新页页顶时会吃掉 \cs{topskip} 基线，注释表首行比正文页首低一行}
% 表头盒要清成 void 而不是留个空盒。\pkg{longtable} 没写 \cs{endhead} 时表头
% 是零高空盒，仍会作为盒子进主竖列；断页恰好断在它前面时，新页第一个盒子
% 就是它，\cs{topskip} 的基线给了它，真正的首行再往下一个行距。清成 void
% 后 \pkg{longtable} 的 \cs{ifvoid} 守卫直接跳过，不再有这只盒子。
%
% \changes{v3.2a}{2026/08/18}{\env{eqdenote} 里局部关掉 \pkg{hyperref} 给
%   \pkg{longtable} 加的锚点：锚点自成一个带撑竿的空段落，表被推到新页时这个
%   空段落占住页首基线，注释首行比正文页首低一行}
% \pkg{longtable} 的 \cs{LT@start} 给每张表排一个链接锚
% （\cs{MakeLinkTarget}，\cs{LTcaptype} 为空时也照排一个
% \texttt{LT@tables} 的），锚变成一个只装撑竿的空段落（\cs{showoutput} 里
% 是 12.2+7.3 的空行盒，内容只有段落撑竿那根零宽 rule，后跟
% \texttt{dest (table.1.1)}）。表整个被推到新页时，空段落抢走 \cs{topskip}
% 基线，「式中」首行低一个行距（1--7 页版实测 140.5 对正文页首 119.9）。
% 注释表不是编号插表，锚没有用处，环境组内把 \cs{MakeLinkTarget} 换成
% 吞参数的空操作，空段落连同锚一起消失。
%    \begin{macrocode}
    \box_gclear:N \LT@head
    \box_gclear:N \LT@firsthead
    \cs_if_exist:NT \MakeLinkTarget
      { \cs_set_eq:NN \MakeLinkTarget \use_none:n }
    \begin{xltabular} { \textwidth }
      { @{} l @{ \hspace { \ccwd } } r @{ \hit@eqdenote@dash } X @{} }
      #1 & #2 \\
    \end{xltabular}
  }
%    \end{macrocode}
%
% 每条注释一个 \cs{item}，不必自己写 \cs{\textbackslash\textbackslash}，也不必记着
% 末条不加。做法是把环境正文按 \cs{item} 切开，再用「\cs{\textbackslash
% \textbackslash} 加一个空格子」接起来。切分只认顶层的 \cs{item}，嵌在花括号里的
% 不算。正文里一个 \cs{item} 都没有时会给出警告，因为那样第二行起缺一个空格子，
% 表格会错位。
%
% 收尾的三句是给后面的元素定位用的。整张表是段落里的一个高盒子，深度取自表格行的
% strut，比正文行深，\hologo{TeX} 只好退到 \cs{lineskip}，末行会与后面的正文挨得比正文
% 行距近；而标题的 \cs{addvspace} 是从盒子底量的，又会偏松。把盒子底与
% \cs{prevdepth} 都收到正文行的深度，接正文和接标题就都跟正文一致。
%    \begin{macrocode}
\NewDocumentEnvironment { eqdenote } { O{\hit@eqdenote@label} +b }
  {
    \seq_set_split:Nnn \l__hit_eqdenote_seq { \item } {#2}
    \seq_get_left:NN \l__hit_eqdenote_seq \l__hit_eqdenote_tl
    \tl_if_blank:VTF \l__hit_eqdenote_tl
      { \seq_pop_left:NN \l__hit_eqdenote_seq \l__hit_eqdenote_tl }
      { \msg_warning:nn { hithesis } { eqdenote-no-item } }
    \tl_set:Nn \l__hit_eqdenote_tl {#1}
    \tl_if_blank:nT {#1}
      { \tl_set:Nn \l__hit_eqdenote_tl { \hphantom { \hit@eqdenote@label } } }
    \exp_args:NV \__hit_eqdenote_table:nn \l__hit_eqdenote_tl
      { \seq_use:Nn \l__hit_eqdenote_seq { \\ & } }
%    \end{macrocode}
%
% \changes{v3.2a}{2026/08/18}{\env{eqdenote} 收尾从「中」字形深度改成段落撑竿
%   深度，后接任何元素的间距与正文相同}
% 收尾要把表底整理成「一行正文的底」。正文行靠段落撑竿把高深撑成正好一个
% \cs{baselineskip}（高 12.22、深 7.30），后面的行距胶、标题的
% \cs{addvspace} 都按这个深度记账。原来按「中」字形深度收，与撑竿体系
% 对不上：接正文 20.9（应 19.45），接小节标题 26.3（应 31.0）。现在量一只
% 段落撑竿的盒子，把 \cs{prevdepth}（此刻等于末行行深）kern 到撑竿深度，
% 再把 \cs{prevdepth} 设成撑竿深度——之后不管接段落、标题还是公式，
% 几何上与接在一段正文后面完全相同。撑竿宏不存在时不动，保持原样。
%    \begin{macrocode}
    \par
    \cs_if_exist:NT \hit@format@strut:
      {
        \hbox_set:Nn \l__hit_eqdenote_box { \hit@format@strut: }
        \dim_gset:Nn \g__hit_eqdenote_dp_dim
          { \box_dp:N \l__hit_eqdenote_box }
        \group_insert_after:N \__hit_eqdenote_finish:
      }
  }
  { }
%    \end{macrocode}
%
% 修正必须推迟到环境组外，而且只设 \cs{prevdepth} 不补 kern。
% \cs{prevdepth} 是组内状态，组一关就被还原，组内设的那份白设；表格行的
% X 列 parbox 里段落撑竿照常生效，行盒的深度本来就是正文行深 7.30，不需要
% kern 找齐（先补过一道 kern，间距反而多出约 7\,bp）。表格收尾后
% \cs{prevdepth} 被 \pkg{longtable} 留成别的值（实测 0.42），所以还是要把
% 撑竿深度存进全局量，\cs{group\_insert\_after:N} 在 \cs{end}\marg{eqdenote}
% 关组之后设回去。
%    \begin{macrocode}
\dim_new:N \g__hit_eqdenote_dp_dim
\cs_new_protected:Npn \__hit_eqdenote_finish:
  { \tex_prevdepth:D \g__hit_eqdenote_dp_dim }
%</shared-math>
%    \end{macrocode}
%
%
% \Finale
